\documentclass[12pt]{article}
\usepackage[a4paper,margin=1in]{geometry}
\usepackage{iftex}
\ifXeTeX
\usepackage{fontspec}
\defaultfontfeatures{Ligatures=TeX}
\setmainfont{Times New Roman}
\else
\usepackage[T1]{fontenc}
\usepackage[utf8]{inputenc}
\usepackage{newtxtext}
\usepackage{newtxmath}
\fi
\usepackage{enumitem}
\usepackage{hyperref}
\usepackage{parskip}
\usepackage{microtype}
\usepackage[normalem]{ulem}
\usepackage{xcolor}

\newcommand{\practiceid}[1]{\texttt{#1}}
\newlength{\readinggutter}
\setlength{\readinggutter}{2.8em}
\newlength{\readingfirstindent}
\setlength{\readingfirstindent}{1.5em}
\newlength{\questionblockgap}
\setlength{\questionblockgap}{0.45\baselineskip}
\newlength{\exampleblockgap}
\setlength{\exampleblockgap}{0.65\baselineskip}
\newlength{\passageparagraphgap}
\setlength{\passageparagraphgap}{0.45\baselineskip}
\newenvironment{exampleblock}{\par\vspace{0.35\baselineskip}\begingroup\raggedright\setlength{\parskip}{0pt}\setlength{\parindent}{0pt}}{\par\endgroup\vspace{\exampleblockgap}}
\newcommand{\readinglineindent}[2]{%
\par\noindent%
\hangindent=\dimexpr\readinggutter+0.9em\relax%
\hangafter=1%
\makebox[\readinggutter][r]{\scriptsize\color{gray}#1}\hspace{0.9em}\hspace*{\readingfirstindent}#2%
\par}
\newcommand{\readingline}[2]{%
\par\noindent%
\hangindent=\dimexpr\readinggutter+0.9em\relax%
\hangafter=1%
\makebox[\readinggutter][r]{\scriptsize\color{gray}#1}\hspace{0.9em}#2%
\par}

\setlength{\parindent}{0pt}
\setlength{\parskip}{0.5em}
\setlist[enumerate]{topsep=0.25em,itemsep=0.18em,parsep=0pt,partopsep=0pt}
\setlength{\emergencystretch}{2em}

\begin{document}
\section*{TOEFL Practice 02 - Tryout}
Source of truth: DOCX only.
\section*{Listening Comprehension}
\subsection*{Part A}
DIRECTIONS: In Part A you will hear short conversations between two speakers. At the end of each conversation, a third speaker will ask a question about what the first two speakers said. Each conversation and each question will be spoken only one time. Therefore, you must listen carefully to understand what each speaker says. After you hear a conversation and the question, read the four choices and select the one that is the best answer to the question the speaker asked. Then, on your answer. sheet, find the number of the question and blacken the space that corresponds to the letter for the answer you have chosen. Blacken the space completely so that the letter inside the space does not show.\par
\begin{exampleblock}
Listen to the following example.\par
On the recording, you hear:\par
(Man) Does the car need to be filled?\par
(Woman) Mary stopped at the gas station on her way home.\par
(Narrator) What does the woman mean?\par
In your test book, you will read:\par
(A) Mary bought some food.\par
(B) Mary had car trouble.\par
(C) Mary went shopping.\par
(D) Mary bought some gas.\par
From the conversation you learn that Mary stopped at the gas station on her way home. The best answer to the question "Does the car need to be filled?" is (D), "Mary bought some gas." Therefore, the correct answer is (D).\par
Now let us begin Part A with question number 1.\par
\end{exampleblock}
\begin{exampleblock}
\textbf{Transkrip rekonstruksi AI (draft; bukan resmi)}\par
Man: Did Peter finally pay Bill?\par
Woman: No. He still hasn't even paid the phone company.\par
Narrator: What does the woman imply about Peter?\par
\end{exampleblock}
\noindent\textbf{1.}
\begin{enumerate}[label=(\Alph*),leftmargin=*,itemsep=0.2em]
\item Peter didn't pay Bill two times.
\item Peter didn't pay the phone company.
\item Peter is the second in line to pay.
\item Peter forgot to call Bill.
\end{enumerate}
\par
\vspace{0.25\baselineskip}
\begin{exampleblock}
\textbf{Transkrip rekonstruksi AI (draft; bukan resmi)}\par
Woman: I tried to register for Professor Allen's course, but the office said I needed permission.\par
Man: That's right. The instructor decides who can enroll in that class.\par
Narrator: What does the man mean?\par
\end{exampleblock}
\noindent\textbf{2.}
\begin{enumerate}[label=(\Alph*),leftmargin=*,itemsep=0.2em]
\item The course is closed for registration.
\item The instructor decides who can enroll.
\item Registration for this course is permitted.
\item The instructor doesn't give students permission.
\end{enumerate}
\par
\vspace{0.25\baselineskip}
\begin{exampleblock}
\textbf{Transkrip rekonstruksi AI (draft; bukan resmi)}\par
Man: How many people came to the workshop?\par
Woman: We expected only fourteen, but almost forty people showed up.\par
Narrator: What does the woman mean?\par
\end{exampleblock}
\noindent\textbf{3.}
\begin{enumerate}[label=(\Alph*),leftmargin=*,itemsep=0.2em]
\item We expected only fourteen people to come.
\item Forty people came to work in the shop.
\item More people came than had been expected.
\item We expanded the shop to include forty people.
\end{enumerate}
\par
\vspace{0.25\baselineskip}
\begin{exampleblock}
\textbf{Transkrip rekonstruksi AI (draft; bukan resmi)}\par
Woman: Did Betty decide what to do about the job offer?\par
Man: Yes. After thinking it over, she decided to take the job.\par
Narrator: What does the man say about Betty?\par
\end{exampleblock}
\noindent\textbf{4.}
\begin{enumerate}[label=(\Alph*),leftmargin=*,itemsep=0.2em]
\item It's better to change this job.
\item Betty decided to take the job.
\item In my opinion, Betty should change her job.
\item I thought Betty was taking John with her.
\end{enumerate}
\par
\vspace{0.25\baselineskip}
\begin{exampleblock}
\textbf{Transkrip rekonstruksi AI (draft; bukan resmi)}\par
Man: Did you call the cleaners?\par
Woman: Yes, I asked whether they had cleaned Jack's jacket yet.\par
Narrator: What did the woman ask?\par
\end{exampleblock}
\noindent\textbf{5.}
\begin{enumerate}[label=(\Alph*),leftmargin=*,itemsep=0.2em]
\item If she dropped Jack off.
\item If she took his jacket to be cleaned.
\item If she was leaving Jack.
\item If she had cleaned his jacket yet.
\end{enumerate}
\par
\vspace{0.25\baselineskip}
\begin{exampleblock}
\textbf{Transkrip rekonstruksi AI (draft; bukan resmi)}\par
Man: Was the teachers' meeting on Wednesday crowded?\par
Woman: Not at all. Most of the teachers never received the notice, so they didn't go.\par
Narrator: What does the woman imply?\par
\end{exampleblock}
\noindent\textbf{6.}
\begin{enumerate}[label=(\Alph*),leftmargin=*,itemsep=0.2em]
\item The teachers didn't go to the meeting.
\item The meeting on Wednesday was crowded.
\item The teachers knew about the meeting.
\item Everyone forgot about the meeting.
\end{enumerate}
\par
\vspace{0.25\baselineskip}
\begin{exampleblock}
\textbf{Transkrip rekonstruksi AI (draft; bukan resmi)}\par
Man: I heard Paul did something nice for you last weekend.\par
Woman: He did. He gave a party for me after my presentation.\par
Narrator: What does the woman say Paul did?\par
\end{exampleblock}
\noindent\textbf{7.}
\begin{enumerate}[label=(\Alph*),leftmargin=*,itemsep=0.2em]
\item He was nice to give up his part.
\item He gave a party for the man.
\item He was kind to do his part.
\item He was kind to come to the party.
\end{enumerate}
\par
\vspace{0.25\baselineskip}
\begin{exampleblock}
\textbf{Transkrip rekonstruksi AI (draft; bukan resmi)}\par
Woman: Could you pick up paper plates at the drug store?\par
Man: I don't think so. Drug stores don't usually sell paper goods like that.\par
Narrator: What does the man mean?\par
\end{exampleblock}
\noindent\textbf{8.}
\begin{enumerate}[label=(\Alph*),leftmargin=*,itemsep=0.2em]
\item Pharmacists shouldn't carry boxes to the station.
\item Paper goods are not usually sold in drug stores.
\item Mobile pharmacies are not stationed here.
\item The pharmacy needs to order paper goods.
\end{enumerate}
\par
\vspace{0.25\baselineskip}
\begin{exampleblock}
\textbf{Transkrip rekonstruksi AI (draft; bukan resmi)}\par
Woman: How was the exam?\par
Man: It looked simple at first, but it was much harder than I expected.\par
Narrator: What does the man mean?\par
\end{exampleblock}
\noindent\textbf{9.}
\begin{enumerate}[label=(\Alph*),leftmargin=*,itemsep=0.2em]
\item The exam seems to be easier than it first appeared.
\item He thinks the exam appears to cover the material.
\item He thinks they made the exam difficult on purpose.
\item The exam is more difficult than he thought.
\end{enumerate}
\par
\vspace{0.25\baselineskip}
\begin{exampleblock}
\textbf{Transkrip rekonstruksi AI (draft; bukan resmi)}\par
Man: What did the new labor report say?\par
Woman: It said that the number of unemployed people has been increasing.\par
Narrator: What does the woman mean?\par
\end{exampleblock}
\noindent\textbf{10.}
\begin{enumerate}[label=(\Alph*),leftmargin=*,itemsep=0.2em]
\item The number of the unemployed has been increasing.
\item Drama clubs have been helping the unemployed.
\item Unemployed actors can find work in drama.
\item Dramas about the unemployed are gaining popularity.
\end{enumerate}
\par
\vspace{0.25\baselineskip}
\begin{exampleblock}
\textbf{Transkrip rekonstruksi AI (draft; bukan resmi)}\par
Woman: Linda seems satisfied with her grades this semester.\par
Man: Really? To me, her grades leave a lot to be desired.\par
Narrator: What does the man imply about Linda's grades?\par
\end{exampleblock}
\noindent\textbf{11.}
\begin{enumerate}[label=(\Alph*),leftmargin=*,itemsep=0.2em]
\item Linda's grades don't leave much room for improvement.
\item Linda deserves better grades in her courses.
\item Linda's grades aren't as high as they should be.
\item Linda's grades should be left alone.
\end{enumerate}
\par
\vspace{0.25\baselineskip}
\begin{exampleblock}
\textbf{Transkrip rekonstruksi AI (draft; bukan resmi)}\par
Man: How was your hamburger?\par
Woman: It was so salty that now I need something to drink.\par
Narrator: What does the woman mean?\par
\end{exampleblock}
\noindent\textbf{12.}
\begin{enumerate}[label=(\Alph*),leftmargin=*,itemsep=0.2em]
\item She had to add salt to her hamburger.
\item She had a hamburger at 1:30.
\item She needs something to drink.
\item Her hamburger was excellent.
\end{enumerate}
\par
\vspace{0.25\baselineskip}
\begin{exampleblock}
\textbf{Transkrip rekonstruksi AI (draft; bukan resmi)}\par
Woman: I forgot to bring my lunch today.\par
Man: Don't worry. When we go out, we'll bring you back a sandwich.\par
Narrator: What do the speakers say they will do?\par
\end{exampleblock}
\noindent\textbf{13.}
\begin{enumerate}[label=(\Alph*),leftmargin=*,itemsep=0.2em]
\item They dropped her sandwich.
\item She didn't have food on her mind.
\item They will bring her a sandwich.
\item She should be on her way home.
\end{enumerate}
\par
\vspace{0.25\baselineskip}
\begin{exampleblock}
\textbf{Transkrip rekonstruksi AI (draft; bukan resmi)}\par
Woman: Is Doug going to vote in the election?\par
Man: No. He won't be eighteen until next year.\par
Narrator: What does the man mean about Doug?\par
\end{exampleblock}
\noindent\textbf{14.}
\begin{enumerate}[label=(\Alph*),leftmargin=*,itemsep=0.2em]
\item Doug is too old to go on a boat.
\item Ducks can't be transported by boat.
\item Doug doesn't like to vote.
\item Doug is too young to vote.
\end{enumerate}
\par
\vspace{0.25\baselineskip}
\begin{exampleblock}
\textbf{Transkrip rekonstruksi AI (draft; bukan resmi)}\par
Man: Is Harry working full-time now?\par
Woman: No, he still spends most of his time at school.\par
Narrator: What does the woman imply about Harry?\par
\end{exampleblock}
\noindent\textbf{15.}
\begin{enumerate}[label=(\Alph*),leftmargin=*,itemsep=0.2em]
\item Harry is working to buy new shoes.
\item Harry's shoes are at school.
\item Harry is a student.
\item Harry walks to the seashore.
\end{enumerate}
\par
\vspace{0.25\baselineskip}
\begin{exampleblock}
\textbf{Transkrip rekonstruksi AI (draft; bukan resmi)}\par
Woman: Is the meeting in Room 307?\par
Man: No. It's in Room 7, but that's on the second floor, not the third.\par
Narrator: What does the man mean?\par
\end{exampleblock}
\noindent\textbf{16.}
\begin{enumerate}[label=(\Alph*),leftmargin=*,itemsep=0.2em]
\item There is no room on the third floor.
\item Room 7 is ready for the third meeting.
\item The room is on a different floor.
\item The floor in this room is dirty.
\end{enumerate}
\par
\vspace{0.25\baselineskip}
\begin{exampleblock}
\textbf{Transkrip rekonstruksi AI (draft; bukan resmi)}\par
Woman: Are those the glasses you always wear?\par
Man: No. I broke my regular pair, so I've been using these instead.\par
Narrator: What does the man mean?\par
\end{exampleblock}
\noindent\textbf{17.}
\begin{enumerate}[label=(\Alph*),leftmargin=*,itemsep=0.2em]
\item After he broke a glass, he had to use a cup.
\item After the flat, he's been using a spare tire.
\item When he broke the glass, he cut himself.
\item He has been using different glasses.
\end{enumerate}
\par
\vspace{0.25\baselineskip}
\begin{exampleblock}
\textbf{Transkrip rekonstruksi AI (draft; bukan resmi)}\par
Man: Have you decided what to get your father for his birthday?\par
Woman: Not exactly, but my brother came up with a good idea for a gift.\par
Narrator: What does the woman mean?\par
\end{exampleblock}
\noindent\textbf{18.}
\begin{enumerate}[label=(\Alph*),leftmargin=*,itemsep=0.2em]
\item My brother had an idea for our father's gift.
\item My brother wants to buy our father a pool.
\item We wanted to bake a pie for my father.
\item We celebrated our father's birthday.
\end{enumerate}
\par
\vspace{0.25\baselineskip}
\begin{exampleblock}
\textbf{Transkrip rekonstruksi AI (draft; bukan resmi)}\par
Woman: Does anyone live in the apartment next to Ann?\par
Man: Not anymore. The people who lived there moved out last week.\par
Narrator: What does the man say about Ann?\par
\end{exampleblock}
\noindent\textbf{19.}
\begin{enumerate}[label=(\Alph*),leftmargin=*,itemsep=0.2em]
\item Ann will be 21 next month.
\item No one now lives next to Ann.
\item Ann moved here three weeks ago.
\item Ann is going on a 21-day vacation.
\end{enumerate}
\par
\vspace{0.25\baselineskip}
\begin{exampleblock}
\textbf{Transkrip rekonstruksi AI (draft; bukan resmi)}\par
Man: I got grease all over this shirt while I was fixing my bike.\par
Woman: Would you like me to sell you something that removes grease?\par
Narrator: What did the woman ask the man?\par
\end{exampleblock}
\noindent\textbf{20.}
\begin{enumerate}[label=(\Alph*),leftmargin=*,itemsep=0.2em]
\item If he could sell her a product to remove grease.
\item If he knew he had a stain.
\item If he wanted to remove his shirt.
\item If she had a crease in her shirt.
\end{enumerate}
\par
\vspace{0.25\baselineskip}
\begin{exampleblock}
\textbf{Transkrip rekonstruksi AI (draft; bukan resmi)}\par
Man: I cleaned the house, and I thought we could play tennis later.\par
Woman: That sounds fine to me.\par
Narrator: What does the woman mean?\par
\end{exampleblock}
\noindent\textbf{21.}
\begin{enumerate}[label=(\Alph*),leftmargin=*,itemsep=0.2em]
\item She approves.
\item He is good at tennis.
\item He was nice to clean the house.
\item She wants to come, too.
\end{enumerate}
\par
\vspace{0.25\baselineskip}
\begin{exampleblock}
\textbf{Transkrip rekonstruksi AI (draft; bukan resmi)}\par
Woman: Are you going out of town this weekend?\par
Man: No. Some friends from Boston are coming to stay with me.\par
Narrator: What does the man mean?\par
\end{exampleblock}
\noindent\textbf{22.}
\begin{enumerate}[label=(\Alph*),leftmargin=*,itemsep=0.2em]
\item He is going out of town.
\item He is expecting guests.
\item He is changing companies.
\item He is moving to another town.
\end{enumerate}
\par
\vspace{0.25\baselineskip}
\begin{exampleblock}
\textbf{Transkrip rekonstruksi AI (draft; bukan resmi)}\par
Man: Could we see the menu, please?\par
Woman: Of course. I'll bring it right over, and today's specials are on the board.\par
Narrator: Where does this conversation probably take place?\par
\end{exampleblock}
\noindent\textbf{23.}
\begin{enumerate}[label=(\Alph*),leftmargin=*,itemsep=0.2em]
\item At an amusement park.
\item In a restaurant.
\item On the sidewalk.
\item Near a parking meter.
\end{enumerate}
\par
\vspace{0.25\baselineskip}
\begin{exampleblock}
\textbf{Transkrip rekonstruksi AI (draft; bukan resmi)}\par
Woman: Should we put the tomatoes next to the beans?\par
Man: Yes, but we should turn the soil before we plant anything.\par
Narrator: What are the speakers probably discussing?\par
\end{exampleblock}
\noindent\textbf{24.}
\begin{enumerate}[label=(\Alph*),leftmargin=*,itemsep=0.2em]
\item Buying vegetables.
\item Planting a garden.
\item Cooking a meal.
\item Loading a truck.
\end{enumerate}
\par
\vspace{0.25\baselineskip}
\begin{exampleblock}
\textbf{Transkrip rekonstruksi AI (draft; bukan resmi)}\par
Man: This area is a little noisy.\par
Woman: Then where would you prefer to sit?\par
Narrator: What does the woman ask the man?\par
\end{exampleblock}
\noindent\textbf{25.}
\begin{enumerate}[label=(\Alph*),leftmargin=*,itemsep=0.2em]
\item Can you hear this noise?
\item Do you know this sound?
\item Where would you prefer?
\item Where is this sound coming from?
\end{enumerate}
\par
\vspace{0.25\baselineskip}
\begin{exampleblock}
\textbf{Transkrip rekonstruksi AI (draft; bukan resmi)}\par
Woman: Excuse me, where can I find the new paperback novels?\par
Man: They're on the second shelf, right next to the dictionaries.\par
Narrator: Where does this conversation probably take place?\par
\end{exampleblock}
\noindent\textbf{26.}
\begin{enumerate}[label=(\Alph*),leftmargin=*,itemsep=0.2em]
\item At a craft show
\item At a bookstore
\item At a hardware store
\item At a video rental shop
\end{enumerate}
\par
\vspace{0.25\baselineskip}
\begin{exampleblock}
\textbf{Transkrip rekonstruksi AI (draft; bukan resmi)}\par
Man: Could you fill it up and check the oil, please?\par
Woman: Sure. I'll also take a look at the tire pressure.\par
Narrator: Who is the man probably speaking to?\par
\end{exampleblock}
\noindent\textbf{27.}
\begin{enumerate}[label=(\Alph*),leftmargin=*,itemsep=0.2em]
\item A gas station attendant
\item A university course grader
\item An income tax accountant
\item A technical team leader
\end{enumerate}
\par
\vspace{0.25\baselineskip}
\begin{exampleblock}
\textbf{Transkrip rekonstruksi AI (draft; bukan resmi)}\par
Man: I can't believe you forgot to send Mom a birthday card.\par
Woman: Well, you knew it was her birthday too. You could have sent one yourself.\par
Narrator: What does the woman mean?\par
\end{exampleblock}
\noindent\textbf{28.}
\begin{enumerate}[label=(\Alph*),leftmargin=*,itemsep=0.2em]
\item She is upset that she forgot to send the card.
\item She thinks the man should share responsibility.
\item She has been forgetting a lot of things lately.
\item She is asking whether the man sent his own card.
\end{enumerate}
\par
\vspace{0.25\baselineskip}
\begin{exampleblock}
\textbf{Transkrip rekonstruksi AI (draft; bukan resmi)}\par
Man: Do you want to go out for dinner?\par
Woman: I'd like to, but I can't leave until I finish my report.\par
Narrator: What will the woman probably do first?\par
\end{exampleblock}
\noindent\textbf{29.}
\begin{enumerate}[label=(\Alph*),leftmargin=*,itemsep=0.2em]
\item Go for dinner
\item Order Italian food
\item Mail her report
\item Finish her report
\end{enumerate}
\par
\vspace{0.25\baselineskip}
\begin{exampleblock}
\textbf{Transkrip rekonstruksi AI (draft; bukan resmi)}\par
Man: I heard Susan gave two weeks' notice.\par
Woman: Yes. She has accepted a position with another company.\par
Narrator: What does the woman imply?\par
\end{exampleblock}
\noindent\textbf{30.}
\begin{enumerate}[label=(\Alph*),leftmargin=*,itemsep=0.2em]
\item She believes her services should be noticed.
\item She is taking a two-week leave.
\item She was promoted to the position of a manager.
\item She is planning to take another job.
\end{enumerate}
\par
\vspace{0.25\baselineskip}
\subsection*{Part B}
DIRECTIONS: In this part of the test, you will hear longer conversations. After each conversation, you will hear several questions. The conversations and questions will not be repeated. After you hear a question, read the four possible answers in. your test book and select. The best answer. Then, on your answer sheet, find the number of the question and fill in the space that corresponds to the letter of the answer you have chosen. Remember, you are not allowed to take notes or write in your test book.\par
\begin{exampleblock}
Listen to the following example:\par
You will hear:\par
You will read:\par
(A) He has changed jobs.\par
(B) He has two children.\par
(C) He has two jobs.\par
(D) He is looking for a job.\par
From the conversation you learn that Tom has taken an additional job. The best answer to the question "Why is Tom tired?" is (C), "He has two jobs." Therefore the correct answer is (C).\par
\end{exampleblock}
\begin{exampleblock}
\textbf{Transkrip rekonstruksi AI (draft; bukan resmi)}\par
Woman: The animated short starts right after the previews.\par
Man: Cartoons again? I came to see the feature film, not a cartoon.\par
Woman: It is part of the program. Besides, modern animation is more complicated than you think.\par
Woman: Why don't you want to watch the animated short?\par
Man: I'm not a child anymore. I'd rather wait for the feature.\par
Woman: You may be surprised. Animated films are not only for children.\par
Woman: Before any drawings are made, the writers create the story.\par
Man: So the plot comes before the artists begin?\par
Woman: Exactly. After the story is set, the scenes are planned and the artists begin their work.\par
Man: When do they add the voices and music?\par
Woman: Usually after the main scenes are prepared. Sound effects, voices, and music are added to make the pictures more lively and to strengthen what the audience sees and hears.\par
\end{exampleblock}
\noindent\textbf{31.}
\begin{enumerate}[label=(\Alph*),leftmargin=*,itemsep=0.2em]
\item In a film studio
\item In a record company
\item At an art gallery
\item At a movie theater
\end{enumerate}
\par
\vspace{0.25\baselineskip}
\noindent\textbf{32.}
\begin{enumerate}[label=(\Alph*),leftmargin=*,itemsep=0.2em]
\item He thinks his children should watch educational programs.
\item His daughters are too old to watch cartoons.
\item He isn't a child any longer.
\item He doesn't want to watch cartoons.
\end{enumerate}
\par
\vspace{0.25\baselineskip}
\noindent\textbf{33.}
\begin{enumerate}[label=(\Alph*),leftmargin=*,itemsep=0.2em]
\item The drawings are made.
\item The story is created.
\item The shots are framed.
\item The artists are organized.
\end{enumerate}
\par
\vspace{0.25\baselineskip}
\noindent\textbf{34.}
\begin{enumerate}[label=(\Alph*),leftmargin=*,itemsep=0.2em]
\item To make cartoons a form of graphic art
\item To enhance the visual and auditory elements
\item To add action to animation
\item To speed up the feature plot
\end{enumerate}
\par
\vspace{0.25\baselineskip}
\begin{exampleblock}
\textbf{Transkrip rekonstruksi AI (draft; bukan resmi)}\par
Woman: Look at all these red and yellow leaves around the trail.\par
Man: This is the best part of walking through the woods in autumn.\par
Woman: It really is beautiful here.\par
Man: The trees only look like this for a few weeks each year.\par
Woman: You mean when the leaves change color before they fall?\par
Man: Right, that's one reason I like this season.\par
Woman: Why do the leaves turn red and yellow?\par
Man: As fall begins, there are fewer hours of daylight. With less light, the green chlorophyll breaks down, and the other colors become visible.\par
Man: So the color change is really a chemical process in the leaves.\par
Woman: Maybe so, but I don't need the whole scientific explanation. I just like the way they look.\par
\end{exampleblock}
\noindent\textbf{35.}
\begin{enumerate}[label=(\Alph*),leftmargin=*,itemsep=0.2em]
\item At a county fair
\item In the woods
\item On a sidewalk.
\item On a playground
\end{enumerate}
\par
\vspace{0.25\baselineskip}
\noindent\textbf{36.}
\begin{enumerate}[label=(\Alph*),leftmargin=*,itemsep=0.2em]
\item In the evening
\item In the morning
\item In the fall
\item In the spring
\end{enumerate}
\par
\vspace{0.25\baselineskip}
\noindent\textbf{37.}
\begin{enumerate}[label=(\Alph*),leftmargin=*,itemsep=0.2em]
\item They shed their leaves in the fall.
\item There are fewer daylight hours.
\item Young trees require a blanket.
\item Red and yellow are prettier than green.
\end{enumerate}
\par
\vspace{0.25\baselineskip}
\noindent\textbf{38.}
\begin{enumerate}[label=(\Alph*),leftmargin=*,itemsep=0.2em]
\item She doesn't care about his explanation.
\item She doesn't understand scientific facts.
\item She doesn't believe the man.
\item She doesn't like autumn.
\end{enumerate}
\par
\vspace{0.25\baselineskip}
\subsection*{Part C}
DIRECTIONS: In Part C you will hear short lectures and extended conversations. At the end of each, you will be asked several questions. Each lecture or conversation and each question will be spoken only one time. For this reason, you must listen carefully to understand what each speaker says. After you hear a question, read the four choices and select the one that best answers the question the speaker asked. Then, on your answer sheet, find the number of the question and blacken the space that contains the letter for the answer you have chosen. Answer all questions according to what is stated or implied in the lecture or conversation.\par
\begin{exampleblock}
Listen to this sample talk.\par
You will hear:\par
Now listen, to the following example.\par
You will hear:\par
You will read:\par
(A) By cars and carriages\par
(B) By bicycles, trains, and carriages\par
(C) On foot and by boat\par
(D) On board ships and trains\par
The best answer to the question "According to the speaker, how did people travel before the invention of the automobile?" is (B), "By bicycles, trains, and carriages." Therefore, the correct answer is (B). You are not allowed to make notes during the test.\par
\end{exampleblock}
\begin{exampleblock}
\textbf{Transkrip rekonstruksi AI (draft; bukan resmi)}\par
Speaker: Many people think air conditioning simply means cooling a room. In fact, its broader purpose is to control the indoor environment so that people can work and live comfortably.\par
Speaker: An air-conditioning system controls not only temperature but also humidity. It can remove moisture from the air when the air is too damp, and it can add moisture when the air is too dry.\par
Speaker: Ventilation is also important. Stale air is taken out of a room and replaced with fresh air from outside.\par
Speaker: In winter, the same basic principles are used. The system still controls temperature, moisture, and ventilation, although the temperature adjustment may involve heating rather than cooling.\par
Speaker: When the indoor environment is properly controlled, workers are less affected by heat, stale air, or excessive humidity. As a result, they often feel more alert.\par
\end{exampleblock}
\noindent\textbf{39.}
\begin{enumerate}[label=(\Alph*),leftmargin=*,itemsep=0.2em]
\item It changes climatic conditions.
\item It controls the indoor environment.
\item It circulates water particles in the air.
\item It reduces overhead expenditures.
\end{enumerate}
\par
\vspace{0.25\baselineskip}
\noindent\textbf{40.}
\begin{enumerate}[label=(\Alph*),leftmargin=*,itemsep=0.2em]
\item By cleaning off dirt and dust
\item By modifying the temperature
\item By removing and adding moisture
\item By bringing in and taking out air
\end{enumerate}
\par
\vspace{0.25\baselineskip}
\noindent\textbf{41.}
\begin{enumerate}[label=(\Alph*),leftmargin=*,itemsep=0.2em]
\item To replace it with fresh air
\item To create a breeze
\item To retard heat conduction
\item To remove local moisture
\end{enumerate}
\par
\vspace{0.25\baselineskip}
\noindent\textbf{42.}
\begin{enumerate}[label=(\Alph*),leftmargin=*,itemsep=0.2em]
\item Hardly ever in cold regions
\item Exclusively in business settings
\item Basically as it does in summer
\item Mainly as needed for repairs
\end{enumerate}
\par
\vspace{0.25\baselineskip}
\noindent\textbf{43.}
\begin{enumerate}[label=(\Alph*),leftmargin=*,itemsep=0.2em]
\item They can dress comfortably.
\item They become tired.
\item They feel alert.
\item They stay healthy.
\end{enumerate}
\par
\vspace{0.25\baselineskip}
\begin{exampleblock}
\textbf{Transkrip rekonstruksi AI (draft; bukan resmi)}\par
Speaker: Although the market for general reporters has become difficult, publications still need people who understand specialized fields. Business and finance, for example, are areas where knowledgeable reporters are needed.\par
Speaker: A reporter may need to interview people, check facts, and meet deadlines. But above all, the reporter must be able to write clearly and accurately.\par
Speaker: In recent years, many newspapers have reduced their staffs. As a result, many reporters have been laid off.\par
Speaker: Even though some traditional newsrooms are shrinking, specialized publications are expected to grow as readers seek more detailed information in particular fields.\par
\end{exampleblock}
\noindent\textbf{44.}
\begin{enumerate}[label=(\Alph*),leftmargin=*,itemsep=0.2em]
\item Business and finance
\item Law and science
\item Sports
\item Chief editors
\end{enumerate}
\par
\vspace{0.25\baselineskip}
\noindent\textbf{45.}
\begin{enumerate}[label=(\Alph*),leftmargin=*,itemsep=0.2em]
\item Copy editing
\item Approaching people
\item Writing
\item Typing
\end{enumerate}
\par
\vspace{0.25\baselineskip}
\noindent\textbf{46.}
\begin{enumerate}[label=(\Alph*),leftmargin=*,itemsep=0.2em]
\item Many reports have been filed.
\item Many reporters have been laid off.
\item Many publications have closed.
\item Many are hiring new specialists.
\end{enumerate}
\par
\vspace{0.25\baselineskip}
\noindent\textbf{47.}
\begin{enumerate}[label=(\Alph*),leftmargin=*,itemsep=0.2em]
\item They will stay sluggish.
\item They may continue to decrease.
\item They are likely to expand.
\item They may be difficult to break into.
\end{enumerate}
\par
\vspace{0.25\baselineskip}
\begin{exampleblock}
\textbf{Transkrip rekonstruksi AI (draft; bukan resmi)}\par
Speaker: Groundwater is one of the world's most important natural resources. For many communities, it represents an essential supply of water for homes, farms, and industry.\par
Speaker: Groundwater supplies are threatened by two major problems: increased use, which can lower water levels, and contamination from chemicals and waste.\par
Speaker: Some groundwater is replaced quickly, but deep underground supplies may take centuries to recharge. Once they are used up or polluted, they cannot be restored in only a few years.\par
\end{exampleblock}
\noindent\textbf{48.}
\begin{enumerate}[label=(\Alph*),leftmargin=*,itemsep=0.2em]
\item It represents an essential supply of water.
\item Its sources can be renewed indefinitely.
\item It always provides a usable source of water.
\item Its development is important in the United States.
\end{enumerate}
\par
\vspace{0.25\baselineskip}
\noindent\textbf{49.}
\begin{enumerate}[label=(\Alph*),leftmargin=*,itemsep=0.2em]
\item Increased rain and agricultural irrigation
\item Increased use and contamination
\item Airborne bacteria and atmospheric oxygen
\item Surface runoff and turbulence
\end{enumerate}
\par
\vspace{0.25\baselineskip}
\noindent\textbf{50.}
\begin{enumerate}[label=(\Alph*),leftmargin=*,itemsep=0.2em]
\item Several years
\item Centuries
\item Dozens of years
\item Thousands of years
\end{enumerate}
\par
\vspace{0.25\baselineskip}
\section*{Structure and Written Expression}
\subsection*{Sentence Completion}
This section is designed to test your ability to recognize language structures that are appropriate in standard written English. The questions in this section belong to two types, each of which has special directions.\par
DIRECTIONS: Questions 1-15 are partial sentences. Below each sentence you will see four words or phrases, marked (A), (B), (C), and (D). Select the one word or phrase that best completes the sentence. Then, on your answer sheet, find the number of the question and fill in the space that contains the letter for the answer you have chosen. Fill in the space completely.\par
\begin{exampleblock}
\textbf{EXAMPLE I}\par
\noindent Drying flowers is the best way......them.
\begin{enumerate}[label=(\Alph*),leftmargin=*,itemsep=0.2em]
\item to preserve
\item by preserving
\item preserve
\item preserved
\end{enumerate}
\vspace{0.25\baselineskip}
The sentence should state, "Drying flowers is the best way to preserve them." Therefore, the correct answer is (A).\par
\end{exampleblock}
\begin{exampleblock}
\textbf{EXAMPLE II}\par
\noindent Many American universities......as small, private colleges.
\begin{enumerate}[label=(\Alph*),leftmargin=*,itemsep=0.2em]
\item begun
\item beginning
\item began
\item for the beginning
\end{enumerate}
\vspace{0.25\baselineskip}
The sentence should state, "Many American universities began as small, private colleges." Therefore, the correct answer is (C). After you read the directions, begin work on the questions.\par
\end{exampleblock}
\noindent\textbf{1.} The upper branches of the tallest trees produce more leaves......other branches.
\begin{enumerate}[label=(\Alph*),leftmargin=*,itemsep=0.2em]
\item than do
\item than does
\item than they do
\item than it does
\end{enumerate}
\par
\vspace{0.25\baselineskip}
\noindent\textbf{2.} No one......projections of demographic shifts are reliable and will prove to be valid in the future.
\begin{enumerate}[label=(\Alph*),leftmargin=*,itemsep=0.2em]
\item know how
\item knows whether
\item knows even
\item know who
\end{enumerate}
\par
\vspace{0.25\baselineskip}
\noindent\textbf{3.} Senior executives often receive bonuses when their profit targets are reached or......
\begin{enumerate}[label=(\Alph*),leftmargin=*,itemsep=0.2em]
\item surpass
\item surpasses
\item surpassed
\item surpassing
\end{enumerate}
\par
\vspace{0.25\baselineskip}
\noindent\textbf{4.} Since the 1970s, riding bicycles......in the United States.
\begin{enumerate}[label=(\Alph*),leftmargin=*,itemsep=0.2em]
\item becomes increasingly widespread
\item become increasingly widely spread
\item has become increasingly widespread
\item has increased and becomes spread widely
\end{enumerate}
\par
\vspace{0.25\baselineskip}
\noindent\textbf{5.} Horseradish has extended stems and a large root that is grated......a spicy food sauce.
\begin{enumerate}[label=(\Alph*),leftmargin=*,itemsep=0.2em]
\item to making
\item to make
\item to be made
\item to the making
\end{enumerate}
\par
\vspace{0.25\baselineskip}
\noindent\textbf{6.} Vitamin A maintains the sharpness of human vision......and promotes healthy bones.
\begin{enumerate}[label=(\Alph*),leftmargin=*,itemsep=0.2em]
\item at night
\item of the night
\item for the night
\item nighttime
\end{enumerate}
\par
\vspace{0.25\baselineskip}
\noindent\textbf{7.} ......, often used in children's poetry and rhymes, are a result of words used in ambiguous contexts.
\begin{enumerate}[label=(\Alph*),leftmargin=*,itemsep=0.2em]
\item Humorous and misunderstood
\item Misunderstand humorously
\item Humorous misunderstandings
\item Misunderstanding its humor
\end{enumerate}
\par
\vspace{0.25\baselineskip}
\noindent\textbf{8.} Modern scanning technology enables physicians to identify brain disorders earlier......than in the past.
\begin{enumerate}[label=(\Alph*),leftmargin=*,itemsep=0.2em]
\item and more accurate
\item and more accurately
\item accurate and more
\item accurately and more
\end{enumerate}
\par
\vspace{0.25\baselineskip}
\noindent\textbf{9.} Educational toys and games give children an opportunity to enjoy themselves......
\begin{enumerate}[label=(\Alph*),leftmargin=*,itemsep=0.2em]
\item while their learning
\item while learning
\item are they learning
\item and they are learning
\end{enumerate}
\par
\vspace{0.25\baselineskip}
\noindent\textbf{10.} William Hazlitt's essays, written in......style, appeared between 1821 and 1822.
\begin{enumerate}[label=(\Alph*),leftmargin=*,itemsep=0.2em]
\item vigorously and informally
\item vigorous and informally
\item vigor and informality
\item vigorous and informal
\end{enumerate}
\par
\vspace{0.25\baselineskip}
\noindent\textbf{11.} Wild hogs inhabited Europe and other parts of the world......6 million years ago.
\begin{enumerate}[label=(\Alph*),leftmargin=*,itemsep=0.2em]
\item as long
\item as long as
\item then it was
\item than it was
\end{enumerate}
\par
\vspace{0.25\baselineskip}
\noindent\textbf{12.} Electrically charged particles exert a magnetic force on one another even......not in physical contact.
\begin{enumerate}[label=(\Alph*),leftmargin=*,itemsep=0.2em]
\item if there are
\item they are
\item if they are
\item are they
\end{enumerate}
\par
\vspace{0.25\baselineskip}
\noindent\textbf{13.} F. Scott Fitzgerald's early literary success led to extravagant living and......a large income.
\begin{enumerate}[label=(\Alph*),leftmargin=*,itemsep=0.2em]
\item a need for
\item to need for
\item needed for
\item for he needed
\end{enumerate}
\par
\vspace{0.25\baselineskip}
\noindent\textbf{14.} Water fire extinguishers must never be used for fires that involve......
\begin{enumerate}[label=(\Alph*),leftmargin=*,itemsep=0.2em]
\item electrically equipped
\item equipment, electrically
\item electricity equipped
\item electrical equipment
\end{enumerate}
\par
\vspace{0.25\baselineskip}
\noindent\textbf{15.} By the 1300s, the Spanish had learned that gunpowder could......propel an object with incredible force.
\begin{enumerate}[label=(\Alph*),leftmargin=*,itemsep=0.2em]
\item use to
\item be used to
\item been used to
\item using it to
\end{enumerate}
\par
\vspace{0.25\baselineskip}
\subsection*{Error Identification}
DIRECTIONS: In questions 16-40 every sentence has four words or phrases that are underlined. The four underlined portions of each sentence are marked (A), (B), (C), and (D). Identify the one word or phrase that makes the sentence incorrect. Then, on your answer sheet, find the number of the question and fill in the space that contains the letter for the answer you have chosen. Fill in the space completely.\par
\begin{exampleblock}
\textbf{EXAMPLE I}\par
Christopher Columbus \uline{has sailed}(A) from Europe in 1492 and \uline{discovered a}(B) \uline{new land}(C) he \uline{thought to}(D) be India.\par
The sentence should state, "Christopher Columbus sailed from Europe in 1492 and discovered a new land he thought to be India." Therefore, you should choose answer (A).\par
\end{exampleblock}
\begin{exampleblock}
\textbf{EXAMPLE II}\par
\uline{As the roles}(A) of people in society change, \uline{so does}(B) the \uline{rules of}(C) conduct in certain \uline{situations}(D).\par
The sentence should state, "As the roles of people in society change, so do the rules of\par
conduct in certain situations." Therefore, you should choose answer (B).\par
\end{exampleblock}
\noindent\textbf{16.} During a recession, manufacturers may \uline{be forced}(A) to \uline{decrease}(B) the number of their workers \uline{to reduction}(C) their \uline{expenditures}(D).
\par
\vspace{\questionblockgap}
\noindent\textbf{17.} \uline{Tides}(A) constitute a \uline{change in}(B) the level of water in \uline{the oceans}(C) and are caused by the gravitational interaction between heavenly \uline{body}(D).
\par
\vspace{\questionblockgap}
\noindent\textbf{18.} Most people \uline{are surprising}(A) to see \uline{how rapidly}(B) bacteria \uline{can multiply}(C) \uline{under favorable}(D) conditions.
\par
\vspace{\questionblockgap}
\noindent\textbf{19.} Without \uline{water}(A), food, shelter, and \uline{clothing}(B), \uline{person}(C) could not survive a prolonged exposure to the \uline{elements}(D).
\par
\vspace{\questionblockgap}
\noindent\textbf{20.} Ancestor \uline{worship}(A) reflects \uline{a family's}(B) reverence for the \uline{advice}(C) and guidance of its \uline{died}(D) members.
\par
\vspace{\questionblockgap}
\noindent\textbf{21.} \uline{A little}(A) land animals live in the polar \uline{regions}(B) which \uline{are covered}(C) with snow \uline{year round}(D).
\par
\vspace{\questionblockgap}
\noindent\textbf{22.} A dolphin, often \uline{called}(A) a "porpoise," is \uline{considered}(B) to be one of the \uline{bright}(C) among \uline{animals}(D).
\par
\vspace{\questionblockgap}
\noindent\textbf{23.} In Arizona, \uline{regular}(A) \uline{annually}(B) events include \uline{horse}(C) shows, art fairs, and folk \uline{dances}(D).
\par
\vspace{\questionblockgap}
\noindent\textbf{24.} \uline{Artificial}(A) intelligence is concerned with \uline{designing}(B) computer systems that perform such tasks as \uline{reason}(C) and \uline{learning}(D) new skills.
\par
\vspace{\questionblockgap}
\noindent\textbf{25.} Saccharin, \uline{made}(A) from toluene, is \uline{about}(B) 350 \uline{times}(C) as \uline{sweeter}(D) as sugar.
\par
\vspace{\questionblockgap}
\noindent\textbf{26.} Work \uline{affect}(A) intellectual development and personal \uline{characteristics}(B), and \uline{personality}(C) and life events \uline{affect}(D) work.
\par
\vspace{\questionblockgap}
\noindent\textbf{27.} \uline{Because of}(A) attitudes \uline{shape}(B) behavior, \uline{psychologists}(C) want to find out how opinions \uline{are formed}(D).
\par
\vspace{\questionblockgap}
\noindent\textbf{28.} Although Connecticut occupies a \uline{small}(A) area, \uline{its}(B) weather can \uline{vary}(C) from one area to \uline{others}(D).
\par
\vspace{\questionblockgap}
\noindent\textbf{29.} The Great Depression \uline{serves as}(A) an example \uline{of drama}(B) fluctuations in \uline{the balanced}(C) \uline{wage rate}(D).
\par
\vspace{\questionblockgap}
\noindent\textbf{30.} In some states, it \uline{has}(A) against \uline{the law}(B) \uline{to ride}(C) a motorcycle \uline{without a}(D) helmet.
\par
\vspace{\questionblockgap}
\noindent\textbf{31.} The government \uline{provides}(A) financial support \uline{for people}(B) who \uline{are unable}(C) to support \uline{themself}(D).
\par
\vspace{\questionblockgap}
\noindent\textbf{32.} Robert Merton \uline{studied}(A) how \uline{does}(B) society \uline{influences}(C) the \uline{development}(D) of science.
\par
\vspace{\questionblockgap}
\noindent\textbf{33.} The \uline{available}(A) of credit \uline{influences}(B) the rate of economic growth and the \uline{increase}(C) \uline{in prices}(D).
\par
\vspace{\questionblockgap}
\noindent\textbf{34.} Pasteurization is the process \uline{of heating}(A) milk \uline{to destroy}(B) \uline{disease-caused}(C) organisms and \uline{bacteria}(D).
\par
\vspace{\questionblockgap}
\noindent\textbf{35.} Paul Claudel, who \uline{written}(A) books about his \uline{personal}(B) \uline{feelings}(C), was a \uline{leading}(D) French author of the early 1900s.
\par
\vspace{\questionblockgap}
\noindent\textbf{36.} Astronomers do not have sufficient \uline{information}(A) \uline{to determine}(B) \uline{what}(C) the solar system \uline{was created}(D).
\par
\vspace{\questionblockgap}
\noindent\textbf{37.} The soil and climate in the tropics are not \uline{suit}(A) to \uline{produce}(B) large \uline{quantities}(C) of \uline{grain}(D).
\par
\vspace{\questionblockgap}
\noindent\textbf{38.} Some \uline{species of}(A) bats \uline{are dormant}(B) \uline{each days}(C) and active \uline{every night}(D).
\par
\vspace{\questionblockgap}
\noindent\textbf{39.} Paper was \uline{so expensive}(A) during the Middle Ages that \uline{it has}(B) to \uline{be used}(C) \uline{sparingly}(D).
\par
\vspace{\questionblockgap}
\noindent\textbf{40.} Additives are \uline{chemicals}(A) that are \uline{infused into}(B) substances to \uline{preventing}(C) them from \uline{spoiling}(D).
\par
\vspace{\questionblockgap}
\section*{Reading Comprehension}
\begin{center}
\textbf{Time 55 minutes}
\end{center}
DIRECTIONS: In this section you will read several passages. Each passage is followed by a series of questions. For questions 1-50, you need to select the best answer, (A), (B), (C), or (D), to each question. Then, on your answer sheet, find the number of the question and fill in the space that contains the letter of the answer you have selected. Fill in the space completely. Answer all questions following a passage on the basis of what is stated or implied in the passage.\par
\begin{exampleblock}
\small
\sloppy
Read the following passage:\par
A tomahawk is a small ax used as a tool and a weapon by the North American Indian\par
tribes. An average tomahawk was not very long and did not weigh a great deal. Originally,\par
the head of the tomahawk was made of a shaped stone or an animal bone and was mounted on\par
Line a wooden handle. After the arrival of the European settlers, the Indians began to use toma-\par
hawks with iron heads. Indian males and females of all ages used tomahawks to chop and cut wood, pound stakes into the ground to put up wigwams, and do many other chores. Indian warriors relied on tomahawks as weapons and even threw them at their enemies. Some types of tomahawks were used in religious ceremonies. Contemporary American idioms reflect\par
this aspect of American heritage.\par
\end{exampleblock}
\begin{exampleblock}
\textbf{EXAMPLE I}\par
\noindent Early tomahawk heads were made of
\begin{enumerate}[label=(\Alph*),leftmargin=*,itemsep=0.2em]
\item stone or bone
\item wood or sticks
\item European iron
\item religious weapons
\end{enumerate}
\vspace{0.25\baselineskip}
According to the passage, early tomahawk heads were made of stone or bone. Therefore, the correct answer is (A).\par
\end{exampleblock}
\begin{exampleblock}
\textbf{EXAMPLE II}\par
\noindent How has the Indian use of tomahawks affected American daily life today?
\begin{enumerate}[label=(\Alph*),leftmargin=*,itemsep=0.2em]
\item Tomahawks are still used as weapons.
\item Tomahawks are used as tools for.certain jobs.
\item Contemporary language refers to tomahawks.
\item Indian tribes cherish tomahawks as heirlooms.
\end{enumerate}
\vspace{0.25\baselineskip}
The passage states, "Contemporary American idioms reflect this aspect of American heritage." The correct answer is (C). After you read the directions, begin work on the questions.\par
\end{exampleblock}
\subsection*{Questions 1-10}
\begingroup
\small
\sloppy
\setlength{\parskip}{0pt}
\setlength{\parindent}{0pt}
The Globe Theater, where most of Shakespeare's plays were staged and performed, was lo-\par
cated in London. Cuthbert and Richard Burbage built the theater in 1599 with materials\par
left over from the construction of London's first playhouse, the Theater. They constructed\par
the Globe on the south side of the Thames River in the little town of Southwark and counted\par
on making the theater a draw for the locals. Little is known about the architectural design of\par
the theater except what can be deduced from maps and the layout of the plays presented\par
there. It appears that the Globe was either round or polygonal on the outside but most likely\par
round on the inside. In keeping with the contemporary imitations of Roman government\par
buildings, its roof was most probably shaped as a crude dome. It can be further deduced that\par
the structure was decorated with pediments, arches, columns, and ornate staircases with\par
carvings of shells, feathers, and cupids. The size of its audience is projected at as many as\par
3,000 spectators both in the amphitheater and in the balcony. The Globe burned down in\par
1613; it was rebuilt on the same foundation a year later, but its external walls were curved at\par
an angle different from that of the original. The theater was built hastily, and evidently\par
safety was not a top priority for either the engineer or the company. After several nearly fatal\par
accidents, the Globe was tom down for good in 1644.\par
\endgroup
\medskip
\noindent\textbf{1.} This passage most likely came from a longer work on
\begin{enumerate}[label=(\Alph*),leftmargin=*,itemsep=0.2em]
\item English deductive trivia
\item English monumental constructions
\item the history of the English theater
\item notable English disasters
\end{enumerate}
\par
\vspace{0.25\baselineskip}
\noindent\textbf{2.} According to the passage, the Globe Theater was built
\begin{enumerate}[label=(\Alph*),leftmargin=*,itemsep=0.2em]
\item from available contemporary materials
\item from materials remaining from another project
\item on a foundation designed to meet a temporary need
\item with ornaments intended to fool the spectators
\end{enumerate}
\par
\vspace{0.25\baselineskip}
\noindent\textbf{3.} In line 5, the phrase "a draw" is closest in meaning to
\begin{enumerate}[label=(\Alph*),leftmargin=*,itemsep=0.2em]
\item an option
\item an attraction
\item a drawing
\item a donation
\end{enumerate}
\par
\vspace{0.25\baselineskip}
\noindent\textbf{4.} It can be inferred from the passage that the Globe's exact architectural design
\begin{enumerate}[label=(\Alph*),leftmargin=*,itemsep=0.2em]
\item should be reconstituted
\item should be obliterated
\item cannot be determined
\item cannot be disregarded
\end{enumerate}
\par
\vspace{0.25\baselineskip}
\noindent\textbf{5.} In line 8, the word "imitations" is closest in meaning to
\begin{enumerate}[label=(\Alph*),leftmargin=*,itemsep=0.2em]
\item enumeration
\item elimination
\item elaborations
\item emulation
\end{enumerate}
\par
\vspace{0.25\baselineskip}
\noindent\textbf{6.} In line 11, the word "projected" is closest in meaning to
\begin{enumerate}[label=(\Alph*),leftmargin=*,itemsep=0.2em]
\item calculated
\item confirmed
\item embellished
\item entrenched
\end{enumerate}
\par
\vspace{0.25\baselineskip}
\noindent\textbf{7.} The passage suggests that, for its time, the Globe Theater was
\begin{enumerate}[label=(\Alph*),leftmargin=*,itemsep=0.2em]
\item humble
\item harsh
\item austere
\item large
\end{enumerate}
\par
\vspace{0.25\baselineskip}
\noindent\textbf{8.} According to the passage, in how many buildings was the Globe Theater housed during its operation?
\begin{enumerate}[label=(\Alph*),leftmargin=*,itemsep=0.2em]
\item One
\item Two
\item Three
\item Four
\end{enumerate}
\par
\vspace{0.25\baselineskip}
\noindent\textbf{9.} The author implies that the last building housing the Globe was
\begin{enumerate}[label=(\Alph*),leftmargin=*,itemsep=0.2em]
\item dignified
\item unmistakable
\item hazardous
\item haunted
\end{enumerate}
\par
\vspace{0.25\baselineskip}
\noindent\textbf{10.} With which of the following statements is the author most likely to agree?
\begin{enumerate}[label=(\Alph*),leftmargin=*,itemsep=0.2em]
\item The architectural design of the theater was exemplary in the 1600s.
\item The builders did not invest a great deal of thought into the theater design.
\item The theater audience enjoyed plays, as well as the building design.
\item The theater location contributed to the opulence of its design and decorations.
\end{enumerate}
\par
\vspace{0.25\baselineskip}
\subsection*{Questions 11-22}
\begingroup
\small
\sloppy
\setlength{\parskip}{0pt}
\setlength{\parindent}{0pt}
Vitamins, taken in tiny doses, are a major group of organic compounds that regulate the\par
mechanisms by which the body converts food into energy. They should not be confused\par
 with minerals, which are inorganic in their makeup. Although in general the naming of vi-\par
tamins followed the alphabetical order of their identification, the nomenclature of individ-\par
ual substances may appear to be somewhat random and disorganized. Among the 13 vita-\par
mins known today, five are produced in the body. Because the body produces sufficient\par
quantities of some but not all vitamins, they must be supplemented in the daily diet. Al-\par
though each vitamin has its specific designation and cannot be replaced by another com-\par
pound, a lack of one vitamin can interfere with the processing of another. When a lack of\par
even one vitamin in a diet is continual, a vitamin deficiency may result.\par
The best way for an individual to ensure a necessary supply of vitamins is to maintain a\par
balanced diet that includes a variety of foods and provides adequate quantities of all the com-\par
pounds. Some people take vitamin supplements, predominantly in the form of tablets. The\par
vitamins in such supplements are equivalent to those in food, but an adult who maintains a\par
balanced diet does not need a daily supplement. The ingestion of supplements is recom-\par
mended only to correct an existing deficiency due to unbalanced diet, to provide vitamins\par
known to be lacking in a restricted diet, or to act as a therapeutic measure in medical treat-\par
ment. Specifically, caution must be exercised with fat-soluble substances, such as vitamins\par
A and D, because, taken in, gigantic doses, they may present a serious health hazard over a pe-\par
riod of time.\par
\endgroup
\medskip
\noindent\textbf{11.} In line 1, the word "regulate" is closest in meaning to
\begin{enumerate}[label=(\Alph*),leftmargin=*,itemsep=0.2em]
\item control
\item refine
\item refresh
\item confine
\end{enumerate}
\par
\vspace{0.25\baselineskip}
\noindent\textbf{12.} According to the passage, vitamins are
\begin{enumerate}[label=(\Alph*),leftmargin=*,itemsep=0.2em]
\item food particles
\item essential nutrients
\item miscellaneous substances
\item major food groups
\end{enumerate}
\par
\vspace{0.25\baselineskip}
\noindent\textbf{13.} In line 4, the word "nomenclature" is closest in meaning to
\begin{enumerate}[label=(\Alph*),leftmargin=*,itemsep=0.2em]
\item conservation
\item classification
\item concentration
\item clarification
\end{enumerate}
\par
\vspace{0.25\baselineskip}
\noindent\textbf{14.} How many vitamins must be derived from nourishment?
\begin{enumerate}[label=(\Alph*),leftmargin=*,itemsep=0.2em]
\item 5
\item 7
\item 8
\item 13
\end{enumerate}
\par
\vspace{0.25\baselineskip}
\noindent\textbf{15.} The author implies that foods
\begin{enumerate}[label=(\Alph*),leftmargin=*,itemsep=0.2em]
\item supply some but not all necessary vitamins
\item should be fortified with all vitamins
\item are equivalent in vitamin content
\item supplement some but not all necessary vitamins
\end{enumerate}
\par
\vspace{0.25\baselineskip}
\noindent\textbf{16.} In line 7, the phrase "daily diet" is closest in meaning to
\begin{enumerate}[label=(\Alph*),leftmargin=*,itemsep=0.2em]
\item weight loss or gain.
\item sufficient quantities
\item nourishment intake
\item vitamin tablets
\end{enumerate}
\par
\vspace{0.25\baselineskip}
\noindent\textbf{17.} A continual lack of one vitamin in a person's diet is
\begin{enumerate}[label=(\Alph*),leftmargin=*,itemsep=0.2em]
\item contagious
\item desirable
\item preposterous
\item dangerous
\end{enumerate}
\par
\vspace{0.25\baselineskip}
\noindent\textbf{18.} With which of the following statements would the author be most likely to agree?
\begin{enumerate}[label=(\Alph*),leftmargin=*,itemsep=0.2em]
\item A varied diet needs to be supplemented with vitamins.
\item An inclusive diet can provide all necessary vitamins.
\item Vitamins cannot be consistently obtained from food.
\item Vitamins should come from capsules in purified form.
\end{enumerate}
\par
\vspace{0.25\baselineskip}
\noindent\textbf{19.} It can be inferred from the passage that vitamin supplements can be advisable
\begin{enumerate}[label=(\Alph*),leftmargin=*,itemsep=0.2em]
\item in special medical cases
\item in most restricted diets
\item after correcting a dietary deficiency
\item before beginning a therapeutic treatment
\end{enumerate}
\par
\vspace{0.25\baselineskip}
\noindent\textbf{20.} In line 17, the phrase "act as" is closest in meaning to
\begin{enumerate}[label=(\Alph*),leftmargin=*,itemsep=0.2em]
\item play the role of
\item pretend to be
\item fight for
\item attest to the fact that
\end{enumerate}
\par
\vspace{0.25\baselineskip}
\noindent\textbf{21.} The author of the passage implies that
\begin{enumerate}[label=(\Alph*),leftmargin=*,itemsep=0.2em]
\item some vitamins are not fat-soluble
\item vitamins can be taken in very small doses
\item most vitamins are water-soluble
\item all vitamins are found in measured doses
\end{enumerate}
\par
\vspace{0.25\baselineskip}
\noindent\textbf{22.} What does the passage mainly discuss?
\begin{enumerate}[label=(\Alph*),leftmargin=*,itemsep=0.2em]
\item Adopting vitamins to control weight
\item The individual's diet for optimum health
\item Vitamin categorization and medical application
\item The place of vitamins in nutrition
\end{enumerate}
\par
\vspace{0.25\baselineskip}
\subsection*{Questions 23-33}
\begingroup
\small
\sloppy
\setlength{\parskip}{0pt}
\setlength{\parindent}{0pt}
When jazz began to lose its reputation as "low-down" music and to gain well-deserved ac-\par
claim among intellectuals, musicians began to feature many instruments previously consid-\par
ered inappropriate for jazz. Whereas before the 1950s, jazz musicians played only eight ba-\par
sic instruments in strict tempo, in this decade they started to improvise on the flute, electric\par
organ, piccolo, accordion, cello, and even bagpipes, with the rhythm section composed for\par
strings or piano. Big bands no longer dominated jazz, and most changes emerged from small\par
combos, such as the Dave Brubeck Quartet and the Gerry Mulligan Quartet. The Gerry\par
Mulligan Quartet proved that a small, modern band could sound complete without a piano;\par
the rhythm section consisted only of a set of drums and a string bass.\par
Jazz continued to move in new directions during the 1960s. Saxophonist and composer\par
Ornette Coleman led a quartet playing "free" jazz that was atonal. Pianist Cecil Taylor also\par
conducted similar experiments with music, and John Coltrane included melodies from In-\par
dia in his compositions. In the 1970s musicians blended jazz and rock music into fusion jazz\par
which combined the melodies and the improvisations of jazz with the rhythmic qualities of\par
rock 'n' roll, with three or five beats to the bar and in other meters. The form of jazz music\par
was greatly affected by electric instruments and electronic implements to intensify, distort,\par
or amplify their sounds. However, the younger musicians of the time felt compelled to in-\par
clude a steady, swinging rhythm which they saw as a permanent and essential element in\par
great jazz\par
\endgroup
\medskip
\noindent\textbf{23.} Which of the following would be the best title for the passage?
\begin{enumerate}[label=(\Alph*),leftmargin=*,itemsep=0.2em]
\item Popular Beats in Classical and Modern Jazz
\item Quintessential Moments in Jazz Music
\item The Achievements of Famous Jazz Musicians
\item The Rising Prestige and Diversity of Jazz
\end{enumerate}
\par
\vspace{0.25\baselineskip}
\noindent\textbf{24.} In line 2, the word "feature" is closest in meaning to
\begin{enumerate}[label=(\Alph*),leftmargin=*,itemsep=0.2em]
\item profess
\item prohibit
\item protest
\item promote
\end{enumerate}
\par
\vspace{0.25\baselineskip}
\noindent\textbf{25.} The paragraph preceding this passage would most likely describe
\begin{enumerate}[label=(\Alph*),leftmargin=*,itemsep=0.2em]
\item instruments used in jazz
\item instrumental pieces in jazz
\item jazz in the 1940s
\item the origins of jazz
\end{enumerate}
\par
\vspace{0.25\baselineskip}
\noindent\textbf{26.} The author of the passage implies that in the 1950s, jazz musicians
\begin{enumerate}[label=(\Alph*),leftmargin=*,itemsep=0.2em]
\item strictly adhered to its traditions and compositions
\item probably continued with its tempo and instrumentation
\item experimented with rhythm and instruments
\item increased the tempo to keep up with the changes
\end{enumerate}
\par
\vspace{0.25\baselineskip}
\noindent\textbf{27.} The author of the passage mentions all of the following EXCEPT
\begin{enumerate}[label=(\Alph*),leftmargin=*,itemsep=0.2em]
\item bagpipes
\item percussion
\item string bass
\item harpsichord
\end{enumerate}
\par
\vspace{0.25\baselineskip}
\noindent\textbf{28.} It can be inferred from the passage that small jazz bands
\begin{enumerate}[label=(\Alph*),leftmargin=*,itemsep=0.2em]
\item were dominated by large orchestras
\item consisted of drums and a string bass
\item were innovative in their music
\item included modern sound systems
\end{enumerate}
\par
\vspace{0.25\baselineskip}
\noindent\textbf{29.} The author believes that the developments in jazz described in the passage
\begin{enumerate}[label=(\Alph*),leftmargin=*,itemsep=0.2em]
\item should be seen as precocious.
\item should be considered influential
\item appear largely suggestive
\item may be perceived as discrete
\end{enumerate}
\par
\vspace{0.25\baselineskip}
\noindent\textbf{30.} The passage implies that representative jazz musicians
\begin{enumerate}[label=(\Alph*),leftmargin=*,itemsep=0.2em]
\item concentrated on melodious combinations of sounds
\item blended improvisations and sheet music together
\item created and modernized sophisticated devices
\item sought novel techniques in form and content
\end{enumerate}
\par
\vspace{0.25\baselineskip}
\noindent\textbf{31.} According to the passage, the changes in jazz music in the 1970s came from
\begin{enumerate}[label=(\Alph*),leftmargin=*,itemsep=0.2em]
\item another harmonious scale
\item another musical trend
\item ambitious aspirations
\item sound amplifications
\end{enumerate}
\par
\vspace{0.25\baselineskip}
\noindent\textbf{32.} In line 17, the word "compelled" is closest in meaning to
\begin{enumerate}[label=(\Alph*),leftmargin=*,itemsep=0.2em]
\item forced
\item challenged
\item obligated
\item censored
\end{enumerate}
\par
\vspace{0.25\baselineskip}
\noindent\textbf{33.} Which of the following best describes the organization of the passage?
\begin{enumerate}[label=(\Alph*),leftmargin=*,itemsep=0.2em]
\item Chronological innovations in jazz music
\item Definitions of diverse jazz styles
\item A classification of prominent jazz musicians
\item Descriptions and examples to illustrate jazz rhythms
\end{enumerate}
\par
\vspace{0.25\baselineskip}
\subsection*{Questions 34-41}
\begingroup
\small
\sloppy
\setlength{\parskip}{0pt}
\setlength{\parindent}{0pt}
The killdeer is a commonly found shore bird that inhabits the area between southern Can-\par
ada and South America. As with all plovers, its soft contour feathers with barbs and barbules\par
impart a sleek appearance to its body while its down feathers insulate it from the winter cold\par
and the summer heat. The male's loud shrill, which seems to say kill-deer, warns other males\par
away from his territory. Ornithologists do not consider the killdeer a true songbird because\par
its throat muscles are not structured to make melodious notes.\par
Killdeers are distinguished by the two black bands that mark their chest and neck.\par
Camouflaged by their protective grayish brown pigment, killdeers build nests that cradle\par
the eggs and the young in shallow depressions in fields and open meadows. Because their\par
nests lie directly on the ground, the young are able to run about as soon as they hatch. Kill-\par
deer are incubatory creatures and brood their own babies. When a predator approaches the\par
nest or the bird's young, the mother tries to distract the intruder by dragging one of her\par
wings as if it were injured. Farmers are particularly fond of killdeers because they feed on in-\par
sects that damage crops. Because whole flocks of killdeers in the wild have vanished due to\par
 overhunting, game laws have been enacted to protect these plovers from poaching.\par
\endgroup
\medskip
\noindent\textbf{34.} It can be inferred from the passage that killdeer usually live
\begin{enumerate}[label=(\Alph*),leftmargin=*,itemsep=0.2em]
\item in the brush
\item in mountains
\item near oceans
\item near cities
\end{enumerate}
\par
\vspace{0.25\baselineskip}
\noindent\textbf{35.} In line 3, the word "impart" is closest in meaning to
\begin{enumerate}[label=(\Alph*),leftmargin=*,itemsep=0.2em]
\item give
\item import
\item link
\item imprint
\end{enumerate}
\par
\vspace{0.25\baselineskip}
\noindent\textbf{36.} What is the reason given for the bird's name?
\begin{enumerate}[label=(\Alph*),leftmargin=*,itemsep=0.2em]
\item It has distinctive bands.
\item It has a peculiar song.
\item It kills young deer.
\item It eats deer fodder.
\end{enumerate}
\par
\vspace{0.25\baselineskip}
\noindent\textbf{37.} In line 6, the word "melodious" is closest in meaning to
\begin{enumerate}[label=(\Alph*),leftmargin=*,itemsep=0.2em]
\item memorable
\item musical
\item mellow
\item marvelous
\end{enumerate}
\par
\vspace{0.25\baselineskip}
\noindent\textbf{38.} It can be inferred from the passage that killdeers are
\begin{enumerate}[label=(\Alph*),leftmargin=*,itemsep=0.2em]
\item inarticulate
\item inconspicuous
\item irreverent
\item irresolute
\end{enumerate}
\par
\vspace{0.25\baselineskip}
\noindent\textbf{39.} How does the mother bird mislead its enemies?
\begin{enumerate}[label=(\Alph*),leftmargin=*,itemsep=0.2em]
\item By pretending to be vulnerable
\item By blending in with the background
\item By building low-lying nests
\item By scaring them away with her cry
\end{enumerate}
\par
\vspace{0.25\baselineskip}
\noindent\textbf{40.} According to the passage, farmers
\begin{enumerate}[label=(\Alph*),leftmargin=*,itemsep=0.2em]
\item form foundations to protect killdeers
\item appreciate the effects of killdeers
\item camouflage killdeer nests and eggs
\item provide killdeers with food and insects
\end{enumerate}
\par
\vspace{0.25\baselineskip}
\noindent\textbf{41.} Which of the following best describes the author's attitude toward the killdeer?
\begin{enumerate}[label=(\Alph*),leftmargin=*,itemsep=0.2em]
\item Menacing
\item Warm
\item Detached
\item Humorous
\end{enumerate}
\par
\vspace{0.25\baselineskip}
\subsection*{Questions 42-50}
\begingroup
\small
\sloppy
\setlength{\parskip}{0pt}
\setlength{\parindent}{0pt}
In 1752, Benjamin Franklin made his textbook experiment with a brass key and a silk\par
kite that he flew in a thunderstorm to prove that lightning and electricity are the same\par
thing. In 1920, a kite-flying championship for families and individuals was held in London.\par
These two seemingly unrelated events underscore the fact that kites can be flown for both\par
pleasure and scientific purposes. For example, in the 1800s weather bureaus flew kites to re-\par
cord temperature and humidity at certain altitudes. On one occasion, ten kites were strung\par
together and flown at a height of four miles to lift men and carry cameras aloft.\par
The kite's ability to fly depends on its construction and the way that its line is attached.\par
The familiar diamond-shaped kite flies when its covered face is aligned against the wind\par
flow. The line attached to the nose of the kite pulls it into the wind, thus creating the neces-\par
sary angle for the lift force. If the kite's construction and the angle of the air stream are cor-\par
rect, the kite will encounter greater pressure against its face and lower pressure against its\par
back. The difference in the pressure creates a lift that causes the kite to rise until it hangs\par
level from its bridle. Its angle against the wind should be sufficiently large or small to create\par
maximum lift to overcome both drag and gravity. The towing point to which the line is at-\par
tached is important because it sets the kite's angle relative to the air flow. Although the kite\par
must be headed up and into the wind with a velocity of 8 to 20 miles per hour, it can main-\par
tain its position through a tail, a rudder, a keel, vents, or tassels.\par
\endgroup
\medskip
\noindent\textbf{42.} What is the main topic of the passage?
\begin{enumerate}[label=(\Alph*),leftmargin=*,itemsep=0.2em]
\item How kites can be utilized
\item Why kites were spurned
\item What parts kites consist of
\item What makes kites stay aloft
\end{enumerate}
\par
\vspace{0.25\baselineskip}
\noindent\textbf{43.} In line 1, the word "textbook" is closest in meaning to
\begin{enumerate}[label=(\Alph*),leftmargin=*,itemsep=0.2em]
\item textual
\item tentative
\item classic
\item outrageous
\end{enumerate}
\par
\vspace{0.25\baselineskip}
\noindent\textbf{44.} In line 4, the word "seemingly" is closest in meaning to
\begin{enumerate}[label=(\Alph*),leftmargin=*,itemsep=0.2em]
\item ostensibly
\item oncoming
\item optimistic
\item opposite
\end{enumerate}
\par
\vspace{0.25\baselineskip}
\noindent\textbf{45.} In line 7, the word "aloft" is closest in meaning to
\begin{enumerate}[label=(\Alph*),leftmargin=*,itemsep=0.2em]
\item in flight
\item in the flood
\item for the analysis
\item for amusement
\end{enumerate}
\par
\vspace{0.25\baselineskip}
\noindent\textbf{46.} According to the passage, the kite flies when its nose is
\begin{enumerate}[label=(\Alph*),leftmargin=*,itemsep=0.2em]
\item pointed away from the ground
\item pointed into the wind flow
\item balanced with the tail
\item aligned parallel to the wind flow
\end{enumerate}
\par
\vspace{0.25\baselineskip}
\noindent\textbf{47.} What is the necessary condition for the kite to fly?
\begin{enumerate}[label=(\Alph*),leftmargin=*,itemsep=0.2em]
\item The kite must be sufficiently strong to withstand great pressure.
\item The kite must be diamond-shaped, and the wind of a certain velocity.
\item The pressure against its back must be lower than the pressure against its face.
\item The pressure of the air flow must be lower than the weight of the kite.
\end{enumerate}
\par
\vspace{0.25\baselineskip}
\noindent\textbf{48.} According to the passage, the line of the kite is important because it
\begin{enumerate}[label=(\Alph*),leftmargin=*,itemsep=0.2em]
\item lifts the kite's cover and frame into the air space
\item contributes to the shape of the kite and extends it
\item determines the angle between the kite and the air flow
\item conveys the direction of the wind and the air flow
\end{enumerate}
\par
\vspace{0.25\baselineskip}
\noindent\textbf{49.} In line 17, the phrase "headed up" is closest in meaning to
\begin{enumerate}[label=(\Alph*),leftmargin=*,itemsep=0.2em]
\item diverted
\item deviated
\item directed
\item drafted
\end{enumerate}
\par
\vspace{0.25\baselineskip}
\noindent\textbf{50.} The paragraph following the passage most would likely discuss
\begin{enumerate}[label=(\Alph*),leftmargin=*,itemsep=0.2em]
\item fiberglass kites flown in competitions
\item the cords and wires needed for kite flying
\item bowed kites curved on their faces
\item elements of kite design and composition
\end{enumerate}
\par
\vspace{0.25\baselineskip}
\end{document}
